\documentclass[10pt,a4paper]{amsart}
\usepackage[margin=19mm]{geometry}
\usepackage{amsmath,amssymb,mathtools}
\usepackage[scr=rsfs,cal=euler]{mathalfa}
\usepackage{xfrac}
\usepackage[unicode,hidelinks,bookmarks=false]{hyperref}
\newcommand{\ie}{i.e.\ }
\newcommand{\cf}{cf.\ }
\newcommand{\resp}{respectively\ }
\newcommand{\Subsection}{\subsection}
\providecommand{\nobreakdash}{\nobreak\hbox{-}}
\setlength{\emergencystretch}{2em}
\multlinegap0pt
\usepackage{bm}
\usepackage{bbm}
\usepackage{dsfont}
\usepackage{xspace}
\usepackage{cancel}
\usepackage{mathtools}
\usepackage{amsthm}
\makeatletter
\@ifundefined{theorem}{\newtheorem{theorem}{Theorem}}{}
\makeatother
\makeatletter
\def\bff{{\boldsymbol f}}
\def\bn{\boldsymbol n}
\def\btau{{\bm \tau}}
\def\bu{{\boldsymbol u}}
\def\bv{{\boldsymbol v}}
\def\bw{{\boldsymbol w}}
\def\bz{\boldsymbol z}
\renewcommand{\div}{\operatorname{div}}
\def\divv{\text{div}}
\expandafter\ifx\csname namedassumption\endcsname\relax
\newenvironment{namedassumption}[1]
%  {\renewcommand{\thistheoremname}{#1}%
%   \begin{genericthm}}
\fi
\expandafter\ifx\csname plan\endcsname\relax
\newenvironment{plan}
{\bigskip\hrule\bigskip\centerline{\bf PLAN}\begin{quote}\tt}
\fi
\def\sig{{\bm \sigma}}
\def\tbu{{\bu}}
\def\tbw{{\bw}}
\def\tp{{p}}
\def\tsig{\sig}
\makeatother
\title[A family of mixed-mixed strongly conservative finite element]{A family of mixed-mixed strongly conservative finite element methods for Biot's model of consolidation}
\author{Qingguo Hong, Johannes Kraus, Maria Lymbery}
\date{Source: arXiv:2607.16418v1 (CC BY 4.0)}
\begin{document}
\maketitle

\noindent\emph{Source and licence.} A family of mixed-mixed strongly conservative finite element methods for Biot's model of consolidation. Version arXiv:2607.16418v1 (CC BY 4.0), distributed under the Creative Commons Attribution 4.0 International licence, as recorded in the arXiv licence metadata for this exact version and shown on the abstract page https://arxiv.org/abs/2607.16418v1. Full rights notice: licenses/arxiv-2607.16418-CC-BY-4.0.txt.\par\smallskip

\noindent\textit{Editorial scope.} Counted unit: the Theorem whose source label is \texttt{thm:near\_best\_approx} in the article (source0.tex lines 1035--1052, source label \texttt{thm:near\_best\_approx}), delivered together with the article's own complete original proof of it (source0.tex lines 1053--1144, one \texttt{proof} environment of 92 lines). No mathematics is written, repaired or re-derived: statement and proof are the article's own text line for line. Required context retained uncounted: eq:Biot\_weak (lines 534--546); eq:X\_h (lines 836--838); eq:Biot\_fem (lines 843--855); inf-sup-dis (lines 922--924). Every definition, display and label of those lines is retained unchanged. Omitted from this package and not redistributed: 1--533, 547--835, 839--842, 856--921, 925--1034, 1145--1335. Nothing inside a retained excerpt is omitted or altered: every retained body is delivered exactly as the article's lines are written, so no omission receipt of any kind accompanies this package. Citations: the retained excerpts carry no citation command at all, so no bibliography entry of the article is transcribed and no reference is redistributed. Reference adaptation: none was needed and none was made; every reference inside the delivered unit resolves inside it, so no unresolved reference or citation is produced. Printed slips kept literal: no printed word, symbol, hypothesis or number of the counted statement or of its proof was altered. Layout and class: the article's own class is \texttt{amsart} with packages including epsf, geometry, listings, algorithmic, algorithm, multirow, enumerate, booktabs, ulem, cancel, accents, graphicx, amscd, color; the wrapper is the lane's pinned \texttt{amsart} editorial wrapper at 10pt on A4, which restates the theorem-like environments the retained text needs and adds the article's own macro definitions that the retained text needs (\texttt{\textbackslash bff}, \texttt{\textbackslash bn}, \texttt{\textbackslash btau}, \texttt{\textbackslash bu}, \texttt{\textbackslash bv}, \texttt{\textbackslash bw}, \texttt{\textbackslash bz}, \texttt{\textbackslash div}, \texttt{\textbackslash divv}, \texttt{\textbackslash sig}, \texttt{\textbackslash tbu}, \texttt{\textbackslash tbw}, \texttt{\textbackslash tp}, \texttt{\textbackslash tsig}), with extra wrapper packages (bm, bbm, dsfont, xspace, cancel, mathtools). The delivered numbering is the unit's own. Assets: no retained span contains a figure environment, a graphics-inclusion command, a PDF-page inclusion, a table or a TikZ picture; no image file is copied into this package, the package declares no asset dependency and no image file is redistributed with it. Licence: version arXiv:2607.16418v1 of this article is distributed under the Creative Commons Attribution 4.0 International licence, as recorded in the arXiv licence metadata for this exact version and shown on the abstract page https://arxiv.org/abs/2607.16418v1. A full rights notice is delivered at licenses/arxiv-2607.16418-CC-BY-4.0.txt. Characters: every retained excerpt is pure ASCII, so the delivered engine, XeLaTeX with its default Latin Modern text font, reports no missing character.
\begin{subequations}\label{eq:Biot_weak}
\begin{align}   
(A \sig, \bm \tau) + (\bu,\divv \bm \tau) = &  \langle \bu_D, 
\bm \tau \bn \rangle_{\Gamma_{\bu,D}}, \label{eq:Biot_weak1}\\
%
(\div \bm \sigma,\bm v) + (p, \div \bv) = & \; -(\bff,\bv),  \label{eq:Biot_weak2}\\  
%
- (R_p^{-1}  \bw, \bz) + (p,\div \bz) = &
\langle p_{D}, \bz \cdot \bn \rangle_{\Gamma_{p,D}}, \label{eq:Biot_weak3} \\ 
%
(\div \bu, q) + (\div \bw, q) + (S p,q) = & (\tilde{g},q) \label{eq:Biot_weak4} 
\end{align}
\end{subequations}

\begin{equation}\label{eq:X_h}
X_h = \Sigma_h^0 \times \bar{V}_h^0 \times P_h
\end{equation}

\begin{subequations}\label{eq:Biot_fem}
\begin{align}   
(A \sig_h, \bm \tau_h) + (\bu_h,\div \tau_h) = &  \langle \bu_D, 
\bm \tau_h \bn \rangle_{\Gamma_{\bu,D}}, \label{eq:Biot_fem1}\\
%
(\div \sig_h,\bv_h) + (p_h, \div \bv_h) = & \; -(\bff,\bv_h),  \label{eq:Biot_fem2}\\  
%
- (R_p^{-1}  \bw_h, \bz_h) + (p_h,\div \bz_h) = &
\langle p_{D}, \bz_h \cdot \bn \rangle_{\Gamma_{p,D}}, \label{eq:Biot_fem3} \\ 
%
(\div \bu_h, q_h) + (\div \bw_h, q_h) + (S p_h,q_h) = & (\tilde{g},q_h) \label{eq:Biot_fem4} 
\end{align}
\end{subequations}

\begin{align}\label{inf-sup-dis}
\sup_{\substack{\bm \tau_h \in \Sigma_h, \\ (\bv_h,\bz_h)\in \bar{V}_h, \\ q_h\in P_h}}\frac{\mathcal A ((\tsig_h,(\tbu_h,\tbw_h),\tp_h),(\bm \tau_h,(\bv_h,\bz_h),q_h))}{\Vert \bm \tau_h \Vert_\Sigma + \Vert (\bv_h,\bz_h)^T \Vert_{\bar{V}} + \Vert q_h \Vert}\ge c_1 \left(\Vert \tsig_h \Vert_\Sigma + \Vert (\tbu_h,\tbw_h)^T \Vert_{\bar{V}} + \Vert \tp_h \Vert\right),
\end{align}

\begin{theorem}\label{thm:near_best_approx}
Let $X:= \Sigma^0 \times \bar{V}^0 \times L^2(\Omega)$ and $X_h$ be defined by~\eqref{eq:X_h}.
Then there exists a constant $C >0$ such that the solutions $(\sig_h,(\bu_h,\bw_h),p_h)^T \in X_h$ of
Problem~\eqref{eq:Biot_fem} and $(\sig,(\bu,\bw),p)^T \in X$ of Problem~\eqref{eq:Biot_weak} satisfy
the estimate
%
\begin{align}\label{eq:near_best_approx}
\Vert \sig - \sig_h \Vert_{H(\div,\Omega;\mathbb{S})} &+ \Vert \bar{\bu} - \bar{\bu}_h \Vert_{\bar{V}} + \Vert p - p_h \Vert  \nonumber  \\
\le &C \inf_{\btau_h \in \Sigma_h,\bar{\bv}_h \in \bar{V}_h, q_h \in Q_h}
  \left(
\Vert \sig - \btau_h \Vert_{H(\div,\Omega;\mathbb{S})} \right. \left.+ \Vert \bar{\bu} - \bar{\bv}_h \Vert_{\bar{V}} + \Vert p - q_h \Vert
\right).
\end{align} 
Furthermore, we have 
\begin{align}\label{eq:near_best_approx_suboptimal}
 \Vert \bar{\bu} - \bar{\bu}_h \Vert_{\bar{V}} +\Vert p - p_h \Vert \le C h^{k-1}\left(h^2|\sig|_k+h|\div\sig|_k+h|\bar{\bu}|_k+|\div\bar{\bu}|_{k-1}+|p|_{k-1}\right).
 \end{align}
\end{theorem}

\begin{proof}
Subtracting~\eqref{eq:Biot_fem} from~\eqref{eq:Biot_weak} results in
%
\begin{subequations}\label{eq:Biot_error}
\begin{align}   
(A (\sig-\sig_h), \bm \tau_h) + ((\bu-\bu_h),\div \bm \tau_h) = &  0, \label{eq:Biot_error1}\\
%
(\div (\sig-\sig_h),\bv_h) + ((p-p_h), \div \bv_h) = & 0,  \label{eq:Biot_error2}\\  
%
- (R_p^{-1}  (\bw-\bw_h), \bz_h) + ((p-p_h),\div \bz_h) = &  0, \label{eq:Biot_error3} \\ 
%
(\div (\bu-\bu_h), q_h) + (\div (\bw-\bw_h), q_h) + (S (p-p_h),q_h) = & 0  \label{eq:Biot_error4} 
\end{align}
\end{subequations}
for all $(\bm \tau_h,(\bv_h,\bz_h),q_h)^T \in X_h$.

Now, let $\sig_I \in \Sigma_h$, $\bu_I \in V_h$, $\bw_I \in V_h$, and $p_I \in P_h$ be arbitrary. Then, in view of~\eqref{eq:Biot_error}, we have
%
\begin{subequations}\label{eq:Biot_error_est}
\begin{align}   
(A (\sig_I-\sig_h), \bm \tau_h) + (\bu_I-\bu_h,\div \bm \tau_h) = &  (A (\sig_I-\sig), \bm \tau_h) + (\bu-\bu_I,\div \bm \tau_h), \label{eq:Biot_error_est1}\\
%
(\div (\sig_I-\sig_h),\bv_h) + (p_I-p_h, \div \bv_h) = & (\div (\sig_I-\sig),\bv_h) + (p-p_I, \div \bv_h),  \label{eq:Biot_error_est2}\\  
%
- (R_p^{-1}  (\bw_I-\bw_h), \bz_h) + (p_I-p_h,\div \bz_h) = &  - (R_p^{-1}  (\bw_I-\bw), \bz_h) + (p-p_I,\div \bz_h), \label{eq:Biot_error_est3} \\ 
%
(\div (\bu_I-\bu_h), q_h) + (\div (\bw_I-\bw_h), q_h) + & (S (p_I-p_h),q_h) 
=  (\div (\bu-\bu_I), q_h) \nonumber \\
+ & (\div (\bw -\bw_I), q_h) + (S (p-p_I),q_h)  \label{eq:Biot_error_est4} 
\end{align}
\end{subequations}
for all $(\bm \tau_h,(\bv_h,\bz_h),q_h)^T \in X_h$.

With the definitions $\bar{\bu}_h :=(\bu_h,\bw_h)$, $\bar{\bv}_h :=(\bv_h,\bz_h)$ and $\bar{\bu}_I :=(\bu_I,\bw_I)$, $\bar{\bv}_I :=(\bv_I,\bz_I)$, from the identity~\eqref{eq:Biot_error_est}, by using the discrete inf-sup condition~\eqref{inf-sup-dis} and Cauchy-Schwarz inequality, we obtain
%
\begin{align*}
\Vert \sig_I - \sig_h \Vert_{H(\div,\Omega;\mathbb{S})} &+ \Vert \bar{\bu}_I - \bar{\bu}_h \Vert_{\bar{V}} +  \Vert p_I - p_h \Vert \\
\le &
c_1^{-1} \left\lvert
\frac{(A (\sig_I-\sig), \bm \tau_h) + (\bu-\bu_I,\div \bm \tau_h)}
{\Vert \btau_h \Vert_{H(\div,\Omega;\mathbb{S})}+ \Vert \bar{\bv}_h \Vert_{\bar{V}} + \Vert q_h \Vert} \right. 
+  \frac{(\div (\sig_I-\sig),\bv_h) + (p-p_I, \div \bv_h)}
{\Vert \btau_h \Vert_{H(\div,\Omega;\mathbb{S})}+ \Vert \bar{\bv}_h \Vert_{\bar{V}} + \Vert q_h \Vert}  \\
+ & \frac{-(R_p^{-1}  (\bw_I-\bw), \bz_h) + (p-p_I,\div \bz_h)}
{\Vert \btau_h \Vert_{H(\div,\Omega;\mathbb{S})}+ \Vert \bar{\bv}_h \Vert_{\bar{V}} + \Vert q_h \Vert}  \\
+ & \left. \frac{(\div (\bu-\bu_I), q_h) + (\div (\bw -\bw_I), q_h) + (S (p-p_I),q_h)}
{\Vert \btau_h \Vert_{H(\div,\Omega;\mathbb{S})}+ \Vert \bar{\bv}_h \Vert_{\bar{V}} + \Vert q_h \Vert} \right\rvert \\
\le &
c_1^{-1}  \left(
\frac{\lvert(A (\sig_I-\sig), \bm \tau_h)\rvert}{\Vert \btau_h \Vert_{H(\div,\Omega;\mathbb{S})}}
+ \frac{\lvert(\bu-\bu_I,\div \bm \tau_h)\rvert}{\Vert \btau_h \Vert_{H(\div,\Omega;\mathbb{S})}} \right. 
+  \frac{\lvert(\div (\sig_I-\sig),\bv_h)\rvert}{\Vert \bar{\bv}_h \Vert_{\bar{V}}}\\
+& \frac{\lvert(p-p_I, \div (\bv_h + \bz_h))\rvert}{\Vert \bar{\bv}_h \Vert_{\bar{V}}} 
+  \frac{\lvert(R_p^{-1}(\bw_I-\bw), \bz_h)\rvert}{\Vert \bar{\bv}_h \Vert_{\bar{V}}} \\
+ & \left. \frac{\lvert(\div (\bu-\bu_I) + \div (\bw-\bw_I), q_h)\rvert}{\Vert q_h \Vert}
+ \frac{\lvert (S (p-p_I),q_h)\rvert}{\Vert q_h \Vert} \right) \\
\le & c_1^{-1}C_0
\left(
\Vert \sig - \sig_I \Vert_{H(\div,\Omega;\mathbb{S})}
+ \Vert \bu - \bu_I \Vert_V 
+ R_p^{-1} \Vert \bw - \bw_I \Vert_V \right. \\
+ &  \left. \Vert \div (\bu-\bu_I) + \div (\bw-\bw_I) \Vert + (1+S) \Vert p-p_I \Vert \right) \\
\le &  c_1^{-1}C_0
 \left(
\Vert \sig - \sig_I \Vert_{H(\div,\Omega;\mathbb{S})}
+ (1+R_p^{-1}) \Vert \bar{\bu} - \bar{\bu}_I \Vert_{\bar{V }}
+ (1+S) \Vert p-p_I \Vert \right).
\end{align*}
%
Finally, by using the triangle inequality together with the previous estimate, we conclude
%
\begin{align*}
\Vert \sig - \sig_h \Vert_{H(\div,\Omega;\mathbb{S})} &+ \Vert \bar{\bu} - \bar{\bu}_h \Vert_{\bar{V}} + \Vert p - p_h \Vert\\
\le &
\Vert \sig - \sig_I \Vert_{H(\div,\Omega;\mathbb{S})} 
+ \Vert \sig_I - \sig_h \Vert_{H(\div,\Omega;\mathbb{S})} 
+  \Vert \bar{\bu} - \bar{\bu}_I \Vert_{\bar{V}} 
+ \Vert \bar{\bu}_I - \bar{\bu}_h \Vert_{\bar{V}} \\
+ & \Vert p - p_I \Vert
+ \Vert p_I - p_h \Vert \\
\le &
C \inf_{\sig_I \in \Sigma_h, \bar{\bu}_I \in \bar{V}_h, p_I \in Q_h} 
\left( \Vert \sig - \sig_I \Vert_{H(\div,\Omega;\mathbb{S})} \right.
+  \left. \Vert \bar{\bu} - \bar{\bu}_I \Vert_{\bar{V}}
+ \Vert p - p_I \Vert \right),
\end{align*} 
%
which proves the assertion~\eqref{eq:near_best_approx} because $\sig_I \in \Sigma_h$, $\bar{\bu}_I \in \bar{V}_h$, and $p_I \in P_h$ were arbitrary.
{
The bound~\eqref{eq:near_best_approx_suboptimal} then follows from standard interpolation error estimates.
}
\end{proof}

\end{document}
